\documentclass[10pt,a4paper]{article}
\usepackage[T1]{fontenc}\usepackage[utf8]{inputenc}
\usepackage{lmodern,amsmath,amssymb,booktabs,array,microtype}
\pdfmapfile{+lm.map}\pdfmapfile{+cm.map}\pdfmapfile{+cmextra.map}\pdfmapfile{+symbols.map}
\pdfinfoomitdate=1\pdftrailerid{}
\usepackage[a4paper,margin=22mm]{geometry}
\usepackage{xcolor,hyperref,fancyhdr}
\definecolor{ink}{HTML}{183A4A}
\hypersetup{colorlinks=true,linkcolor=ink,urlcolor=ink,pdftitle={CST-DUST: Basisunabhaengiger Kopplungsbefund v0.1}}
\pagestyle{fancy}\fancyhf{}\fancyhead[L]{\small X6 | CST-DUST Kopplungsbefund}
\fancyhead[R]{\small v0.1 | 12.11.2026}\fancyfoot[C]{\thepage}
\setlength{\headheight}{14pt}\setlength{\parindent}{0pt}\setlength{\parskip}{5pt}
\setlength{\emergencystretch}{2em}\renewcommand{\arraystretch}{1.15}
\newcommand{\dd}{\mathrm d}\newcommand{\one}{\mathbf1}\newcommand{\D}{\mathcal D}
\newcommand{\A}{\mathcal A}\newcommand{\B}{\mathsf B}
\begin{document}
{\color{ink}\LARGE\bfseries Ein basisunabhängiger Kopplungsbefund}\par
{\large Vollständiger Differenzblock in einem festgelegten Zustandsbereich}\par
\textbf{Ergebnis: Der gewichtete $6\times6$-Differenzblock ist zertifiziert ungleich null.}

Die Aussage gilt für die ursprüngliche endliche Familie aus 33 AA-Historien,
dieselbe Entkopplung auf Wirkungsebene und explizit festgelegte Eingangs- und
Ausgangstesträume. Alle komplexen Matrixelemente werden eingeschlossen.
Die größte Singularzahl misst die maximal erreichbare absolute Testantwort
bei Einheitsnorm in $V$. Ihre Definition hängt nicht von der gewählten Basis ab.

\section*{Die drei getrennten Ergebnisse}
\begin{center}\begin{tabular}{lrr}\toprule
$V$-normierte Blockgröße & Gesamter $6\times6$-Block & Atom$\to$Positiv, $4\times2$\\\midrule
$\sigma_{\max}(K_1-K_0)$ & $0{,}032602301976$ & $0{,}012323446639$\\
$\sigma_{\max}(R_{\rm conn})$ & $0{,}007251598185$ & $0{,}003652388389$\\
$\sigma_{\max}(R_{\rm mean})$ & $0{,}030294483160$ & $0{,}010489512041$\\\bottomrule
\end{tabular}\end{center}
\textbf{Jede Tabellenzahl besitzt eine deterministische Einschließung
$\boldsymbol{\pm10^{-12}}$, einschließlich der Dezimalrundung.}
Es handelt sich weder um statistische Konfidenzintervalle noch um
Wahrscheinlichkeiten. Die Identität $K_1-K_0=R_{\rm conn}+R_{\rm mean}$ gilt auf
Blockebene. Die drei Singularzahlen werden nicht addiert.

Der zusammenhängende Historienanteil ist ebenfalls von null getrennt.
Eine Änderung nur der mittleren Skalarantwort erklärt den gesamten
Differenzblock deshalb nicht.

\section*{Ein einzelner Nulleintrag entscheidet nicht mehr}
Als zusätzliche, vorher festgelegte Kontrolle werden die bereits konstruierten
strikt positiven Historiengewichte verwendet, die den alten Grundpaketkanal
exakt auslöschen. Mit diesen Gewichten ergibt sich dennoch
\[
 \sigma_{\max}(\Delta_6)=0{,}008446681427\pm2\cdot10^{-12},\qquad
 \sigma_{\max}(\Delta_{+\leftarrow0})=0{,}001702265450\pm2\cdot10^{-12}.
\]
Diese Kontrolle ist ein Befund für genau diese Gewichte, keine uniforme
Untergrenze über alle zulässigen Historienmaße.

\section*{Präziser Geltungsbereich}
Berechnet wird $R_F(K_1-K_0)|_E:E\to F'$, also die auf $F$ eingeschränkte
Antidualantwort. Es ist nicht die Norm einer $H\to H$-Evolution und nicht die
volle Operatornorm auf $V\to V'$. Der Differenzblock hat \textbf{exakt Rang zwei};
vier unabhängige Eingangskombinationen bleiben unempfindlich.
Die Aussage lautet somit: Es existieren nachweislich empfindliche Zustände
innerhalb des deklarierten Bereichs. Sie lautet nicht: Jeder Zustand reagiert.

Physikalische CST-Evolution, ein ausgezeichnetes Historienmaß, Unitarität
und Singularitätsauflösung bleiben offen. Die frühere Entscheidung zur
praktischen kompakten Reduktion bleibt HOLD.

\newpage
\section{Räume, Normen und vollständige Gram-Matrizen}
Die Maße und Skalen des Vorgängers bleiben unverändert. In dimensionslosen
Variablen gelten
\[
 \dd\lambda(t)=\tfrac12\one_{[-1,1]}\dd t,\quad
 \dd\nu(r)=\tfrac12\delta_0(\dd r)+\tfrac12\one_{(0,1)}\dd r,\quad
 \dd\mu_A=\tfrac12\prod_{j\in A}\dd\lambda(t_j)\dd\nu(r_j)\dd q_j.
\]
Mit $u=q_1^2+q_2^2$ sind $h_0=\pi^{-1/2}e^{-u/2}$ und
$h_R=(u-1)h_0$. Beide sind in der ungewichteten Feldnorm orthonormal,
gerade und unter Randvertauschung invariant. Die Dichtepakete lauten
\[
 d_0=2\one_{00},\qquad
 d_+=\frac2{\sqrt3}(\one_{10}+\one_{01}+\one_{11}),\qquad
 d_c=\frac2{\sqrt6}(\one_{10}+\one_{01}-2\one_{11}).
\]
Die vollständigen Zustände $\Psi_{a,b}=\sqrt2\,d_a h_b$ enthalten konstante
Uhrpakete. Das kohärente Vorzeichen in $d_c$ ist Teil der Präparation.
Alle sechs Zustände sind in $H=L^2_{\rm inv}(\dd\mu_A)$ orthonormal.

Die vor der neuen Auswertung festgelegten Räume sind
\[
 E_6=F_6=\operatorname{span}\{\Psi_{a,b}:a\in\{0,+,c\},\ b\in\{0,R\}\}.
\]
Für den separat berichteten Mischblock gelten
$E_2=\operatorname{span}\{\Psi_{0,0},\Psi_{0,R}\}$ und
$F_4=\operatorname{span}\{\Psi_{+,0},\Psi_{+,R},\Psi_{c,0},\Psi_{c,R}\}$.

Die geprüfte Topologie ist
\[
 V=L^2_{\rm inv}((1+u)^4\dd\mu_A)\subset H\subset
 V'=L^2_{\rm inv}((1+u)^{-4}\dd\mu_A).
\]
Da $u$ im Grundpaket radial das Maß $e^{-u}\dd u$ trägt, ergibt sich exakt
\[
 G=\begin{pmatrix}65&196\\196&913\end{pmatrix},\quad
 G_{ij}=\int_0^\infty e^{-u}(1+u)^4 f_i(u)f_j(u)\dd u,
 \quad (f_0,f_R)=(1,u-1).
\]
Es gilt $\det G=20929>0$. Die volle Gram-Matrix ist $G_6=I_3\otimes G$.
\textbf{Die Normierung nur der Diagonale wäre falsch:} Der Überlapp $196$
bleibt sonst unbeachtet. Ein exakter Basiswechsel zur $V$-Orthonormalbasis ist
\[
 T=\begin{pmatrix}
 1/\sqrt{65}&-196/\sqrt{65\cdot20929}\\
 0&\sqrt{65/20929}
 \end{pmatrix},\qquad T^*GT=I_2.
\]
Schreibe $U_6=I_3\otimes T$. Für Rohpaarungen
$\B_{ij}=\langle\Psi_i,K\Psi_j\rangle$ ist die auszuwertende Matrix
$\widehat\B=U_F^*\B U_E$. Ihre Spektralnorm ist
\[
 \|R_FK|_E\|=\sup_{\substack{f\in E,\,g\in F\\\|f\|_V=\|g\|_V=1}}
 |\langle g,Kf\rangle|=\sigma_{\max}(\widehat\B).
\]
$F'$ wird mit der von $V$ induzierten Dualnorm versehen. In der
gewichteten Funktionsdarstellung entspricht es dem Antidual-Unterraum
$(1+u)^4F$. Es wird kein neues physikalisches Skalarprodukt postuliert.

\newpage
\section{Wirkung, Historiensumme und analytische Feldmomente}
Es gelten $L=\{0,1\}$, $I=\{2\}$, $R=\{3,4\}$. Alle 33 abgedeckten,
transitiv geschlossenen AA-Ordnungen werden mit $1/33$ summiert.
Unter $S_2(L)\times S_2(R)$ entstehen zehn Orbits; ihre Multiplizitäten sind
$4,4,4,4,4,4,2,4,2,1$. Der Code enumeriert die gelabelte Familie vollständig.

Mit $\beta=4/\sqrt6$, $c_{ij}=(1,-9,16,-8)_{|I(i,j)|}$ gilt
\[
 H_C^{(0)}=-\tfrac12D^T\operatorname{diag}(c_{ij})D,\quad
 M_C=\beta H_C^{(0)},\quad F_4(C)=5-\sum_{i\prec j}c_{ij}.
\]
Die Intervention wird vor jeder Integration definiert:
\[
 \overline M=\tfrac1{33}\sum_CM_C,\quad
 M_{C,\lambda}=\overline M+\lambda(M_C-\overline M),\quad\lambda\in\{0,1\}.
\]
Gravitation, Staub, Maße, Zustände und Gewichte sind für beide Endpunkte
identisch. Die Referenz $\overline M$ wird auch bei der zusätzlichen
Gewichtskontrolle nicht verändert. Die primären Amplituden werden nicht
durch die komplexe Gravitationssumme geteilt.

Ordne die Feldkoordinaten als $(q_0,q_1,q_3,q_4,q_2)$ und setze
$Q_\epsilon=\operatorname{diag}(1,1,1,1,\epsilon)-iM_{C,\lambda}$.
Der gemeinsame Abel-Grenzwert wird analytisch genommen.
Die innere Lücke ist $\alpha\geq\beta>0$. Vier Rand-Gaußpivots und der
anschließende innere Pivot fixieren die Fresnelphase. Die numerische
Berechnung verifiziert die jeweilige Halbebene, bevor sie die Wurzel nimmt.
Eine Hauptwurzel der gesamten Determinante ohne Phasenverfolgung wird vermieden.

Sei $W=Q_0^{-1}$, seien $d_j$ die fünf Eliminationspivots und
\[
 a=\frac{4\pi\sqrt{2\pi}}{\prod_{j=1}^5\sqrt{d_j}},\quad
 x=W_{00}+W_{11}-1,\quad y=W_{22}+W_{33}-1.
\]
Alle vier Feldübergänge sind dann exakt durch Wick-Momente gegeben:
\[
 \A_C=a\begin{pmatrix}
 1&x\\y&xy+2\sum_{i=0}^1\sum_{j=2}^3W_{ij}^{\,2}
 \end{pmatrix}.
\]
Zeilen bezeichnen ausgehende, Spalten eingehende Feldpakete. Die Quadrate
sind komplexe Quadrate, keine Betragsquadrate. Die entkoppelte Antwort
$\A_{\rm ref}$ folgt aus derselben Rechnung mit $\overline M$.

\section{Dichteintegration und vollständiger Block}
Für $\kappa_j=\tfrac12+\tfrac\beta4\sum_{i\prec j}c_{ij}(t_j-t_i)^2$ gilt
$E(k)=\int_0^1e^{-ikr}\dd r=e^{-ik/2}\operatorname{sinc}(k/2)$.
Schreibe $E_j=E(\kappa_j)$. Die vier ausgehenden Maskenintegrale sind
\[
 (d_{00},d_{10},d_{01},d_{11})_C=
 \frac1{16}\int\dd\lambda^5(t)\,(1+E_2)(1,E_3,E_4,E_3E_4).
\]
Der Faktor $1/16=(1/4)\cdot2\cdot4\cdot(1/32)$ enthält beide
Randquotienten, beide $\sqrt2$-Faktoren, beide Dichtepakete und alle fünf
Atom-/Positivgewichte. Kein inneres Dichtestratum wird ausgelassen.

\newpage
Die symmetrischen ausgehenden Dichteantworten sind
\[
 d_C=\left(d_{00},\frac{d_{10}+d_{01}+d_{11}}{\sqrt3},
                    \frac{d_{10}+d_{01}-2d_{11}}{\sqrt6}\right)^T.
\]
Die beiden Eingangspunkte sind minimal: $\kappa_0=\kappa_1=1/2$.
Deshalb ist ihr vollständiger Dichtefaktor die feste Zeile
\[
 r=\left(1,\frac{2e+e^2}{\sqrt3},\frac{2e-2e^2}{\sqrt6}\right),
 \qquad e=E(1/2).
\]
Für $p_C=e^{iF_4(C)}/33$ lautet somit der gesamte $6\times6$-Block
\[
 \B_1=\sum_Cp_C(d_Cr)\otimes\A_C,\qquad
 \B_0=\sum_Cp_C(d_Cr)\otimes\A_{\rm ref}.
\]
Diese Formel berechnet alle 36 Einträge. Die Reduktion auf zwei
Eingangsfeldrichtungen ist eine exakte Faktorisierung, keine
Zustandsabschneidung. Bei der zusätzlichen Kontrolle ersetzt das vorher
festgelegte Orbitgewicht die jeweilige Multiplizität geteilt durch 33.

\section{Der vollständige normierte Differenzblock}
Die Reihenfolge ist $(0,0),(0,R),(+,0),(+,R),(c,0),(c,R)$, wobei die
Feldbasis durch $T$ orthonormiert ist. Die folgenden Tabellen zeigen
Real- und Imaginärteil von $\widehat\Delta_6$; jede angezeigte Komponente
besitzt einschließlich Tabellenrundung Fehler kleiner als $10^{-9}$.
Die Zertifikate verwenden die viel engeren Rohintervalle in JSON.
\[\small \operatorname{Re}\widehat\Delta_6=\begin{pmatrix}
-0.003092382 & 0.004437978 & -0.003935220 & 0.005729178 & -0.000589893 & 0.000819538\\
0.004432897 & -0.006500059 & 0.005721450 & -0.008477074 & 0.000818987 & -0.001171886\\
-0.002723553 & 0.003921984 & -0.002668617 & 0.003978312 & -0.000783643 & 0.001083599\\
0.003899689 & -0.005706588 & 0.003941271 & -0.005924427 & 0.001082218 & -0.001531651\\
-0.000527872 & 0.000767054 & -0.000721143 & 0.001057400 & -0.000084331 & 0.000119394\\
0.000755172 & -0.001123877 & 0.001038184 & -0.001552395 & 0.000118484 & -0.000173905\\
\end{pmatrix}\]
\[\small \operatorname{Im}\widehat\Delta_6=\begin{pmatrix}
0.001847410 & -0.002503902 & 0.004674645 & -0.006472434 & -0.000417378 & 0.000610953\\
-0.002503147 & 0.003512273 & -0.006468411 & 0.009231204 & 0.000610082 & -0.000907413\\
0.003066696 & -0.004171553 & 0.006425383 & -0.008860579 & -0.000250757 & 0.000380946\\
-0.004173889 & 0.005824353 & -0.008851976 & 0.012498954 & 0.000376666 & -0.000574200\\
0.000226153 & -0.000311468 & 0.000654941 & -0.000924190 & -0.000078486 & 0.000115441\\
-0.000311764 & 0.000450745 & -0.000918085 & 0.001345110 & 0.000113237 & -0.000169598\\
\end{pmatrix}\]

Die zweite Singularzahl des Differenzblocks ist
$0{,}000221198937\pm10^{-11}>0$. Die übrigen vier sind exakt null:
Für beide Feldzustände liegen die Dichtekombinationen
$(-r_+,1,0)^T$ und $(-r_c,0,1)^T$ im Eingangskern.

Der $4\times2$-Atom$\to$Positiv-Block besteht aus Zeilen 3 bis 6 und
Spalten 1 bis 2. Seine vollständigen komplexen Einträge sind damit ebenfalls
explizit vorhanden. Die JSON-Dateien enthalten zusätzlich $K_1$, $K_0$,
$R_{\rm conn}$ und $R_{\rm mean}$, jeweils unnormiert und $V$-normiert.

\newpage
\section{Zwei getrennte Ursachenblöcke}
Setze $Z=\sum_Cp_C$, $\D=\sum_Cp_Cd_Cr$ und
$\widetilde\A=\sum_Cp_C\A_C/Z$. $Z$ ist bei beiden getesteten
Gewichtungen zertifiziert von null getrennt. Dann gilt exakt
\begin{align*}
 \B_1-\B_0&=R_{\rm conn}+R_{\rm mean},\\
 R_{\rm conn}&=\sum_Cp_C(d_Cr)\otimes(\A_C-\widetilde\A),\\
 R_{\rm mean}&=\D\otimes(\widetilde\A-\A_{\rm ref}).
\end{align*}
Dies erweitert die frühere skalare Diagnose auf den vollständigen Block.
Die komplexen Gewichte $p_C/Z$ sind keine Wahrscheinlichkeiten.
Der zusammenhängende Rest ist eine Amplitudendiagnose des festgelegten
Ensembles, kein allgemeiner Verschränkungs- oder Kausalitätsnachweis.
Beide Beiträge sind im neuen Zustandsbereich nachweislich ungleich null.

\section{Zertifizierte Singularzahlen statt einer Float-SVD}
Nach $V$-Normierung besitzt jeder Block die Form
$[r_0C\ \ r_+C\ \ r_cC]$, $C\in\mathbb C^{6\times2}$.
Seine beiden nichtverschwindenden Singularzahlen sind exakt diejenigen
von $\|r\|_2C$. Im Atom$\to$Positiv-Block werden nur die vier positiven
Ausgangszeilen von $C$ verwendet; hier ist $r_0=1$.

Für $H=C^*C=\left(\begin{smallmatrix}a&z\\\bar z&d\end{smallmatrix}\right)$
lauten die Eigenwerte
\[
 \lambda_\pm=\frac{a+d\pm\sqrt{(a-d)^2+4|z|^2}}2.
\]
Alle Operationen werden in 192-Bit-Arb/Acb mit nach außen gerundeten
Intervallen ausgeführt. Beide Eigenwerte sind für alle berichteten
Blöcke positiv eingeschlossen; $\sigma_{\max}=\sqrt{\lambda_+}$.
Die Fehlerradien der Integrale werden bereits vor dieser Rechnung
eintragsweise berücksichtigt.

Eine zusätzliche Zertifizierung arbeitet unabhängig von der Rangfaktorisierung
auf den ursprünglichen Gram-Matrizen:
\[
 H_{\rm gen}=\B^*G_F^{-1}\B,\quad
 \frac{v^*H_{\rm gen}v}{v^*G_Ev}>l^2,
 \qquad u^2G_E-H_{\rm gen}>0.
\]
Ein rationaler komplexer Testvektor beweist die untere Grenze $l$.
Positive führende Hauptminoren in Ballarithmetik beweisen die obere
Grenze $u$ per Sylvester-Kriterium. Float64 bestimmt nur Kandidaten,
nicht die Gültigkeit der Ungleichungen. Alle 20 Blockzertifikate bestehen
diese zweite Prüfung.

\section{Warum der Befund basisunabhängig ist}
Für beliebige invertierbare Basiswechsel $P$ am Eingang und $Q$ am
Ausgang gelten
\[
 \B'=Q^*\B P,\qquad G_E'=P^*G_EP,\qquad G_F'=Q^*G_FQ.
\]
Der Quotient in der Supremumsdefinition bleibt durch die bijektive
Umparametrisierung unverändert. Äquivalent unterscheiden sich die
orthonormierten Matrizen nur durch unitäre Faktoren. Ein komplexer,
nichtunitärer Scher-/Skalierungswechsel wird zusätzlich numerisch
mit den transformierten Gram-Matrizen zertifiziert. Eine naive SVD
der Rohmatrix würde dabei von etwa $0{,}7201$ auf $3{,}6290$ wechseln;
der korrekte Wert bleibt $0{,}032602301976\pm10^{-12}$.

\newpage
\section{Vollständiges Fehlerbudget im deklarierten Bereich}
Feld- und Dichteintegrale sind analytisch. Die fünf Uhrintegrale werden
mit einer positiven, normierten 40-Punkt-Gauß-Legendre-Regel je Achse
ausgewertet. Die Ausgangsuhren werden als exakte Umordnung endlicher
Summen kontrahiert. Der Fehler wird nicht aus der Differenz zweier
Quadraturordnungen abgeleitet.

Für eine Achse in der Bernsteinellipse $\mathcal E_\rho$, $\rho=2$,
gelten $a=(\rho+\rho^{-1})/2$, $b=(\rho-\rho^{-1})/2$.
Mit inzidentem Absolutgrad $W_j$ ist die Majorante einer einzelnen
Maskenpaarung $M_j\leq\tfrac18\exp[\tfrac\beta2b(a+1)W_j]$.
Aus den Chebyshevkoeffizienten und der Exaktheit bis Grad $2n-1$ folgt
\[
 e_C\leq\frac12\frac{\rho^{-2n}}{1-\rho^{-1}}
       \sum_{j=0}^4\exp\!\left[\tfrac\beta2b(a+1)W_j\right].
\]
Für die Ausgangsatommaske sind nur die Paare zu Punkt 2 erforderlich;
die entsprechende engere Majorante wird separat berechnet.
Linearkombinationen der Masken erhalten die Summe der absoluten
Koeffizienten mal Maskenfehler. Für einen skalaren Matrixfaktor $A_C$
gilt eintragsweise $e\leq\sum_C|p_C||A_C|e_C$.
Beim Differenzblock wird direkt $A_C-A_{\rm ref}$ eingesetzt.
Danach werden die vollständigen Gram-Transformationen und die
Singularwertformeln auf die eingeschlossenen Einträge angewendet.

\begin{center}\begin{tabular}{p{66mm}p{82mm}}\toprule
Fehleranteil & Behandlung\\\midrule
Uhrquadratur & Analytische Ellipsenschranke pro Historie und Maske\\
Rundung, Gaußphasen, Normierung & Arb/Acb, 192 Bit; Dezimalexporte nach außen\\
Abel-Extrapolation & Null: analytischer Grenzwert\\
Dichte- und Feldquadratur & Null: analytische Integration und Gaußmomente\\
Historienauslassung & Null innerhalb der festgelegten 33er-Familie\\
Zustandsprojektion & Null für die definierten endlichen Paarungsblöcke\\
Antworten außerhalb von $F$ & Nicht durch dieses endliche Blockzertifikat begrenzt\\\bottomrule
\end{tabular}\end{center}

\section{Prüfungen, Reproduktion und Quellen}
\textbf{13 Kontrollgruppen bestanden.} Dazu gehören die unabhängige direkte
fünfdimensionale 6er-Quadratur aller 33 gelabelten Historien, reelle
Schur-Eigenphasen, kartesische Gram-Momente, Wiedergewinnung alter
Matrixelemente, die beiden Blockzerlegungen, nichtunitäre Basiswechsel,
vier exakte Eingangsnulrichtungen und die frühere Auslöschungskontrolle.
Die kleinen Floatkontrollen testen die Implementierung; die
Genauigkeitszusage stammt aus den analytischen Grenzen und Ballrechnungen.

\texttt{python -B verify.py} prüft Hashes und führt die vollständige Rechnung
in einer temporären Kopie erneut aus. Die drei numerischen Ergebnisdateien
müssen bytegleich sein. Der Vertrag wurde vor der neuen Blockauswertung
fixiert und mit SHA-256 gebunden; das ist kein externer Zeitstempel.

Quellbasis: \href{https://github.com/thomasmgfeller/x6-spectral-astrolabium/tree/d7acdb6a2558bbec0b6c2c208742debc55479e75}{X6-Repository, Main d7acdb6a2558},
insbesondere Randzustände, Mischblock, Definitionsbereich und Reduktionsgrenze;
außerdem das aktuelle Prüfpaket Mischungsqualifikation v0.1.
Exakte Dateiherkunft und Hashes: \texttt{provenance.json}.
Arithmetik: \href{https://python-flint.readthedocs.io/en/latest/arb.html}{Python-FLINT 0.9.0, Arb}
und \href{https://python-flint.readthedocs.io/en/latest/acb_mat.html}{Acb-Matrizen}.
Die hier benötigten Norm-, Moment- und Fehlerformeln sind oben hergeleitet.

\textbf{Grenze des Ergebnisses:} Basisunabhängigkeit gilt innerhalb derselben
Räume und derselben Norm. Raumwechsel, Gewichtswechsel, andere Wirkungen,
unendliche CST-Summen und physikalische Übergangswahrscheinlichkeiten
erfordern eigene Nachweise.
\end{document}
